% SGA 3, Expose V, Loop 1 English reconstruction.
% Scope: Theorem 4.1 proof part b), followed by Lemma 4.1.1 and its proof;
% stops immediately before part c), Proof of (i).
% Authority: Polo--Gille Exp5-13oct24.pdf, component-local pp. 12--13
% (combined-reader pp. 274--275). The re-edited page headers are 262--263;
% internal original-page markers in this range advance through 265--266.
% Jacob Reinhold's CC BY 4.0 Markdown is comparison/drafting material only.

\medskip
\phantomsection\label{proof:sga3-V-4.1-b}%
\noindent\textbf{b)} 考虑 \(A_0\) 的两个素理想 \(x\) 和 \(y\).
我们证明等式
\[
  x\cap B=y\cap B
\]
蕴含存在 \(A_1\) 的素理想 \(z\), 使
\[
  x=d_0(z),\qquad y=d_1(z).
\]

事实上, 若此断言不成立, 则对 \(A_1\) 的每个满足
\(\delta_1^{-1}(t)=y\) 的素理想 \(t\),
\(x\) 均不同于 \(\delta_0^{-1}(t)\).
对每个这样的 \(t\), 都有
\[
  \delta_0^{-1}(t)\cap B
  =\delta_1^{-1}(t)\cap B
  =y\cap B=x\cap B.
\]
由 Cohen--Seidenberg 定理 (参见 Bourbaki,
\emph{Alg\`ebre commutative}, V, \S2, 命题~1 的推论~1)
可知, \(x\) 不包含在任何一个理想 \(\delta_0^{-1}(t)\) 中.%
\footnote{\label{note:sga3-V-21}编者注: 以下对原文作了展开;
特别地, 新增了取自 \hyperref[bib:sga3-V-DG70]{[DG70]}, III,
\S2.4, 引理~4.3 的引理~\hyperref[lem:sga3-V-4.1.1]{4.1.1}.}
\(A_1\) 中满足 \(\delta_1^{-1}(t)=y\) 的素理想 \(t\) 至多有 \(n\) 个
(参见\emph{上述文献}, 命题~3).
因此, 由素理想回避引理 (\emph{上述文献}, II, \S1, 命题~3),
存在元素 \(a\in x\), 不属于任何理想 \(\delta_0^{-1}(t)\).

于是 \(\delta_0(a)\) 不属于这些素理想 \(t\) 中的任何一个,
而 \hyperref[lem:sga3-V-4.1.1]{下面的引理} 将蕴含范数
\(N_{\delta_1}(\delta_0(a))\) 不属于 \(y\).
这里的范数通过 \(\delta_1\) 将 \(A_1\) 视为 \(A_0\)-代数来计算;
沿用 \hyperref[proof:sga3-V-4.1-a]{a)} 的记号,
\[
  N_{\delta_1}(\delta_0(a))=\sigma_n.
\]
另一方面,
\[
  (-1)^{n-1}\sigma_n
  =a^n+\sum_{i=1}^{n-1}(-1)^i\sigma_i a^{n-i},
\]
故此范数属于 \(B\cap x=B\cap y\), 矛盾.

\medskip
\phantomsection\label{lem:sga3-V-4.1.1}%
\noindent\textbf{引理 4.1.1.}---
\textit{设 \(A\to A'\) 是交换环同态, 使 \(A'\) 成为秩为 \(n\)
的投射 \(A\)-模. 设 \(p\in\operatorname{Spec}(A)\),
\(q_1,\ldots,q_r\) 是 \(\operatorname{Spec}(A')\) 中位于 \(p\)
上方的点, 并设 \(a\in A'\). 则 \(a\in q_1\cup\cdots\cup q_r\)
当且仅当其范数 \(N(a)\) 属于 \(p\).}

事实上, 将 \(A\) 与 \(A'\) 替换为局部化 \(A_p\) 与 \(A'_p\),
即可化归到 \((A,p)\) 局部且 \(A'\) 半局部的情形, 此时
\[
  \operatorname{Spec}(A')=\{q_1,\ldots,q_r\}.
\]
在此情形中, \(A'\) 是秩为 \(n\) 的自由 \(A\)-模
(参见 Bourbaki, \emph{Alg\`ebre commutative}, II, \S3.2, 命题~5 的推论~2),
且 \(N(a)\) 是自同态
\[
  \ell_a:A'\longrightarrow A',\qquad a'\longmapsto aa'.
\]
的行列式. 因而有等价关系
\[
  N(a)\notin p
  \quad\Longleftrightarrow\quad N(a)\text{ 可逆}
  \quad\Longleftrightarrow\quad \ell_a\text{ 可逆}
  \quad\Longleftrightarrow\quad
  a\notin q_1\cup\cdots\cup q_r.
\]
% Translation statistics: source_lines=86; output_lines=85; inline_math=49; display_math=8; equation=0; equation_star=0; align=0; footnotes=1.
% Review: complete assigned component; pending_TODO=0; all math symbols, labels, references, diagram paths and environments checked; formula prose translated.
